\section{Full Proof of the T2 Lower Bound}\label{app:t2-proof}

This appendix proves \Cref{thm:t2}.  The proof is self-contained: all target
families, auxiliary games, weights, activation rules, deletion errors, and
asymptotic estimates are defined below.

Throughout this appendix $n$ is large, $0<\delta<1/4$ is fixed, and
\[
  Y:=n^\delta.
\]
Let $\mathcal P_Y$ be the set of odd primes at most $Y$.  For
$a<c$ in $\mathcal P_Y$, put
\[
  I_{a,c}:=\left(\frac{n}{2ac},\frac{n}{ac}\right],
  \qquad
  B_{a,c}:=I_{a,c}\cap\PP,
\]
and define the target family
\[
  \mathcal T
  :=
  \{\,acb:\ a<c,\ a,c\in\mathcal P_Y,\ b\in B_{a,c}\,\}.
\]
For large $n$, every such $b$ is greater than $Y$: indeed
$b>n/(2Y^2)=\frac12 n^{1-2\delta}>Y$.  Thus $a,c,b$ are distinct primes and
each target $acb$ lies in $(n/2,n]$.

\begin{proposition}[Finite scaled capture game]\label{prop:t2-finite-capture}
Let $H$ be a finite weighted incidence system.  Its vertices form a finite set
$V$, its edges form a finite set $\mathcal E$, each edge $e$ has a non-negative
weight $w(e)$, and each incidence set $e\subseteq V$ has size at most $3$.
Consider the following Maker-first scored game.

At any time some vertices are captured, some vertices are deleted, and some
edges are scored.  An edge is live if it is not scored and is disjoint from the
deleted vertices.  Maker chooses a live edge $f$, scores $w(f)$, and captures
all vertices of $f$.  Breaker then either deletes one uncaptured vertex or
scores one live edge.  Let $S$ be the total scored weight, by either player,
and for a live edge $e$ define
\[
  \Phi(e):=2^{|e\cap C|}w(e),
\]
where $C$ is the current captured-vertex set.  Scored and deleted edges
contribute no live potential.  Set
\[
  Q:=8S+\sum_{e\text{ live}}\Phi(e).
\]
Maker has a strategy such that after each full Maker--Breaker round $Q$ is at
least its value before Maker's move.  Consequently, if the game is continued
until no live positive-weight edge remains, the final scored weight is at least
\[
  \frac18\sum_{e\in\mathcal E}w(e).
\]
The same conclusion holds for the rank-$2$ graph variant in which Breaker may
delete an unscored live edge instead of scoring it, and Maker-scored edges are
called claimed edges.
\end{proposition}

\begin{proof}
We prove the rank-$3$ scored statement; the graph version is the same argument
with edge deletion replacing Breaker's scored-edge move.  The graph edge
deletion loses exactly the current live potential of that edge, while a
scored-edge move changes $Q$ by $8w(e)-\Phi(e)\ge 0$.

For a live Maker move $f$, let
\[
  \Delta(f):=Q(\text{after Maker scores }f)-Q(\text{before Maker scores }f).
\]
Maker chooses a live edge maximizing $\Delta(f)$.  Such a choice exists because
the game is finite.

We record the two local domination facts used in the round estimate.  First,
if $g$ is live and $|g\cap C|\le 2$, then scoring $g$ itself increases $Q$ by
at least its current live potential $\Phi(g)$.  Indeed, the contribution of
$g$ changes from $2^r w(g)$ to $8w(g)$ with $r=|g\cap C|\le 2$, and
$8w(g)-2^r w(g)\ge 2^r w(g)$.  All other live-edge contributions can only
increase when Maker captures new vertices.

Second, let $v$ be a vertex that Breaker could delete after Maker's move, and
let $\Lambda_v$ be the live potential incident with $v$ at that moment.  If
$\Lambda_v>0$, choose a live edge $g$ containing $v$ after Maker's move.  The
edge $g$ was also live before Maker's move.  Compare Maker's actual max-gain
edge with the alternative move $g$.  When Maker scores $g$, every live edge
through $v$ at the beginning of the round gains its whole current potential
from the capture of $v$, and the selected edge $g$ supplies enough additional
gain to cover the possible coefficient increases caused by the actual move
before $v$ is deleted.  Edge by edge, this gives
\[
  \Delta(g)\ge \Lambda_v.
\]
Since Maker chose an edge of maximum $\Delta$, the actual Maker gain is at
least $\Lambda_v$.

Now consider Breaker's reply.  If Breaker scores a live edge, $Q$ does not
decrease, as observed above.  If Breaker deletes an uncaptured vertex $v$, the
only lost live terms are the live edges incident with $v$, whose total
potential is $\Lambda_v$ after Maker's move.  The previous paragraph shows
that Maker's already-secured gain is at least $\Lambda_v$, so the full
Maker--Breaker round does not decrease $Q$.

Initially $C=\varnothing$, $S=0$, and every edge is live, so
\[
  Q_0=\sum_{e\in\mathcal E}w(e).
\]
At terminal time there is no positive-weight live edge, hence
$Q_{\mathrm{fin}}=8S_{\mathrm{fin}}$.  Monotonicity gives
$8S_{\mathrm{fin}}\ge Q_0$, which is the desired one-eighth bound.
\end{proof}

\begin{proposition}[Initial target mass]\label{prop:t2-initial-mass}
For every fixed $0<\delta<1/4$,
\[
  W_0
  :=
  \sum_{\substack{a<c\\a,c\in\mathcal P_Y}} |B_{a,c}|
  \gg_\delta
  \frac{n(\log\log n)^2}{\log n}.
\]
\end{proposition}

\begin{proof}
Since $ac\le Y^2=n^{2\delta}$, the interval endpoint
$X_{a,c}:=n/(ac)$ satisfies $X_{a,c}\ge n^{1-2\delta}\to\infty$.  By the prime
number theorem, uniformly for all such pairs and all sufficiently large $n$,
\[
  \pi(X_{a,c})-\pi(X_{a,c}/2)
  \ge
  \kappa_\delta \frac{X_{a,c}}{\log n}
  =
  \kappa_\delta \frac{n}{ac\log n}
\]
for some positive constant $\kappa_\delta$.  Therefore
\[
  W_0
  \ge
  \kappa_\delta\frac{n}{\log n}
  \sum_{\substack{a<c\\a,c\in\mathcal P_Y}}\frac1{ac}.
\]
Mertens' theorem for primes gives
\[
  \sum_{p\in\mathcal P_Y}\frac1p
  =
  \log\log Y+O(1)
  =
  \log\log n+O_\delta(1),
\]
and $\sum_p p^{-2}<\infty$.  Hence
\[
  \sum_{a<c\in\mathcal P_Y}\frac1{ac}
  =
  \frac12\left[
    \left(\sum_{p\in\mathcal P_Y}\frac1p\right)^2
    -
    \sum_{p\in\mathcal P_Y}\frac1{p^2}
  \right]
  =
  \left(\frac12+o(1)\right)(\log\log n)^2.
\]
Combining the last two displays proves the claim.
\end{proof}

\begin{proposition}[Residual divisibility embedding]\label{prop:t2-residual-embedding}
Let $\mathcal T_*$ be any subfamily of targets $t=acb\in\mathcal T$ such that,
for every $t=acb\in\mathcal T_*$, the numbers $a$, $c$, and $ac$ are already
unavailable as legal moves.  Attach to $t$ the three slot labels
\[
  b,\qquad ab,\qquad cb
\]
and let $e_t:=\{b,ab,cb\}$.  Then the residual divisibility game on
$\mathcal T_*$ is at least as favorable to Prolonger as the scored rank-$3$
hypergraph game on the hyperedges $\{e_t:t\in\mathcal T_*\}$:
\begin{enumerate}
  \item the only future moves that can make $t=acb$ illegal are
    $b$, $ab$, $cb$, and $t$ itself;
  \item playing $b$, $ab$, or $cb$ deletes exactly the residual targets whose
    hyperedges contain that slot;
  \item playing the exact target $t$ kills no other target in $\mathcal T_*$,
    and is therefore a scored-edge move;
  \item every live hyperedge corresponds to a legal actual move $t$.
\end{enumerate}
\end{proposition}

\begin{proof}
Fix $t=acb\in\mathcal T_*$.  Because $t>n/2$, no proper multiple of $t$ lies in
$\{2,\dots,n\}$.  Thus a future move can make $t$ illegal only by being $t$
itself or a proper divisor of $t$.  The positive divisors are
\[
  1,\ a,\ c,\ ac,\ b,\ ab,\ cb,\ acb.
\]
The divisor $1$ is not a move, while $a$, $c$, and $ac$ are unavailable by
hypothesis.  This proves the first assertion.

For slot incidence, suppose another target is $t'=a'c'b'$.  Since
$b,b'>Y\ge a,a',c,c'$, the divisibility $b\mid t'$ holds iff $b=b'$.  Thus a
move $b$ kills exactly the targets in the $b$-fiber.  Similarly, if
$ab\mid a'c'b'$, the large prime factor forces $b=b'$, and after cancellation
$a$ must be $a'$ or $c'$.  This is exactly the condition that the slot $ab$
belongs to $e_{t'}$.  The proof for $cb$ is identical.

If two targets $t,u$ in $\mathcal T_*$ are comparable, say $t\mid u$, then
$u$ is a multiple of $t$ lying in $(n/2,n]$.  The next positive multiple of
$t$ after $t$ is $2t>n$, so $u=t$.  Hence distinct targets are pairwise
incomparable, and playing an exact target scores only that target.

Finally, if $e_t$ is live in the slot hypergraph, then none of $b$, $ab$, $cb$,
or $t$ has been played.  The only other proper divisors of $t$ are
$1,a,c,ac$, and these are either not moves or already unavailable.  There is no
proper multiple of $t$ on the board.  Therefore no played number is comparable
with $t$, so the actual move $t$ is legal.  This proves the comparison.
\end{proof}

\begin{proposition}[Activation-stage potential]\label{prop:t2-activation-potential}
During activation, let each unclaimed pair edge $e=(a,c)$ carry the current
token set $B_e(t)$ of primes $b\in B_{a,c}$ for which the target $acb$ has not
yet been killed.  Write $w_t(e):=|B_e(t)|$.

Prolonger plays as follows.  In each activation round, among pair edges with
positive current token weight, he chooses a max-gain edge in the graph version
of \Cref{prop:t2-finite-capture}; if $e=(a,c)$ is chosen, he plays one live
target $acb$ with $b\in B_e(t)$.  This scores one actual move, claims the pair
edge $e$, captures the vertices $a$ and $c$, and leaves the remaining live
tokens on that secured pair for the residual phase.

Let $S_t$ be the number of activation targets already played by Prolonger, and
define
\[
  Q_t
  :=
  S_t+\sum_e \phi_t(e),
\]
where
\[
  \phi_t(e)=
  \begin{cases}
    \frac18 w_t(e),& e\text{ unclaimed and neither endpoint captured},\\[1mm]
    \frac14 w_t(e),& e\text{ unclaimed and one endpoint captured},\\[1mm]
    \frac12 w_t(e),& e\text{ unclaimed and two endpoints captured},\\[1mm]
    w_t(e),& e\text{ claimed},\\[1mm]
    0,& e\text{ deleted}.
  \end{cases}
\]
Let $E$ be the total number of target tokens deleted during activation by
Shortener moves not modeled as graph vertex deletions $a$ or graph edge
deletions $ac$.  At the end of activation,
\[
  Q_{\mathrm{end}}\ge \frac{W_0}{8}-E.
\]
Equivalently,
\[
  S_{\mathrm{end}}
  +
  \sum_{e\text{ secured}} w_{\mathrm{end}}(e)
  \ge
  \frac{W_0}{8}-E.
\]
\end{proposition}

\begin{proof}
First ignore off-model token deletions.  The activation process is exactly the
rank-$2$ graph game of \Cref{prop:t2-finite-capture}, with one harmless
bookkeeping change: when Prolonger claims an edge $e$, he consumes one token
from $e$ and adds it immediately to the score $S_t$.  If the coefficient of
$e$ before the move is $\gamma\in\{1/8,1/4,1/2\}$, the contribution of $e$
changes from $\gamma w_t(e)$ to
\[
  1+(w_t(e)-1)=w_t(e),
\]
exactly as if the edge had become claimed with full weight.  All other
incident edges change as in the graph game.  Thus the scaled potential
argument gives $Q_{\mathrm{end}}\ge W_0/8$ when every harmful Shortener move is
a graph deletion.

Now restore all actual Shortener moves.  A token $acb$ is deleted precisely
when Shortener plays one of its still-available harmful divisors or the target
itself.  Deleting one live token can reduce $Q_t$ by at most $1$, since every
coefficient appearing in the definition of $\phi_t(e)$ is at most $1$.
Therefore all off-model deletions together reduce the lower bound by at most
$E$, giving
\[
  Q_{\mathrm{end}}\ge W_0/8-E.
\]
At the end of activation, a positive-weight edge is either secured or deleted,
so $Q_{\mathrm{end}}$ is exactly the activation score plus the residual live
target weight on secured pairs.
\end{proof}

\begin{proposition}[Activation deletion budget]\label{prop:t2-deletion-budget}
For fixed $0<\delta<1/4$, the off-model deletion count from
\Cref{prop:t2-activation-potential} satisfies
\[
  E\ll \frac{Y^4}{\log^4 Y}
  =
  o\!\left(\frac{n(\log\log n)^2}{\log n}\right).
\]
\end{proposition}

\begin{proof}
There is at most one activation round per small-prime pair, so the number of
rounds is
\[
  R\le \binom{|\mathcal P_Y|}{2}
  \ll \frac{Y^2}{\log^2Y}.
\]
We bound the number of live tokens an off-model Shortener move can delete.

If Shortener plays a large prime $b$, then for each pair $(a,c)$ there is at
most one token $acb$ using that $b$.  Hence the move deletes at most
$O(Y^2/\log^2Y)$ tokens.  If Shortener plays a lateral semiprime $pb$ with
$p\le Y$ and $b>Y$ prime, the killed targets have the form $pcb$ with the
third small prime varying, so there are at most $O(Y/\log Y)$ of them.  If
Shortener plays an exact target $acb$, exactly one token is deleted.

The worst of these bounds is the large-prime bound.  Across $R$ activation
rounds,
\[
  E\ll R\frac{Y^2}{\log^2Y}
  \ll \frac{Y^4}{\log^4Y}.
\]
Since $Y=n^\delta$,
\[
  \frac{Y^4/\log^4Y}{n(\log\log n)^2/\log n}
  \ll_\delta
  \frac{n^{4\delta-1}}{(\log n)^3(\log\log n)^2}
  \to 0
\]
because $4\delta-1<0$.
\end{proof}

\begin{proposition}[Residual mass on secured pairs]\label{prop:t2-residual-mass}
At the end of activation, the secured pairs carry residual live target weight
\[
  M:=
  \sum_{e\text{ secured}} w_{\mathrm{end}}(e)
  \gg_\delta
  \frac{n(\log\log n)^2}{\log n}.
\]
\end{proposition}

\begin{proof}
By \Cref{prop:t2-activation-potential},
\[
  S_{\mathrm{end}}+M\ge \frac{W_0}{8}-E.
\]
The activation score is at most the number of small-prime pairs:
\[
  S_{\mathrm{end}}\le R\ll \frac{Y^2}{\log^2Y}
  =
  o\!\left(\frac{n(\log\log n)^2}{\log n}\right),
\]
because $2\delta<1$.  The deletion budget $E$ is negligible by
\Cref{prop:t2-deletion-budget}, while $W_0$ has the required lower bound by
\Cref{prop:t2-initial-mass}.  Subtracting the two negligible terms gives the
claim.
\end{proof}

\begin{proposition}[T2 lower bound]\label{prop:t2-final}
For every fixed $0<\delta<1/4$ there is $c_\delta>0$ such that
\[
  \L(n)\ge c_\delta\frac{n(\log\log n)^2}{\log n}
\]
for all sufficiently large $n$.
\end{proposition}

\begin{proof}
After activation, restrict to the residual live targets lying over secured
pairs.  Their total weight is $M$, as in
\Cref{prop:t2-residual-mass}.  For each secured pair $(a,c)$, the activation
move made $a$, $c$, and $ac$ unavailable.  Therefore
\Cref{prop:t2-residual-embedding} compares the remaining divisibility game to a
scored rank-$3$ slot hypergraph whose total edge weight is $M$.

By \Cref{prop:t2-finite-capture}, Prolonger can force scored residual target
weight at least $M/8$ in this hypergraph model.  The comparison proposition says
that every such scored hypergraph move is a genuine actual game move, and that
the actual game gives Prolonger no fewer legal options than the model.  Hence
the number of residual moves after activation is at least $M/8$.  Since
$M\gg_\delta n(\log\log n)^2/\log n$, the same lower bound holds for the total
game length.

Choosing any fixed $\delta<1/4$, for instance $\delta=1/8$, gives an absolute
constant $c>0$ and proves \Cref{thm:t2}.
\end{proof}
